\section{Что остаётся открытым и как устроен следующий поиск}

\paragraph{Доказанные границы.}
Выпрямление всех гармонических вложений не получено. В частности,
$H(3,7,8)$, $H(3,7,11)$ и $H(4,5,7)$ в этом досье не имеют
ни удостоверенного выпрямителя, ни доказательства невыпрямляемости.
Невозможность одного переноса, наличие критических точек у одного
полинома параметра и неудача ограниченной системы Грёбнера не заменяют
ни одного из этих двух результатов. Для утверждений о неприводимости
случая нужно исключить все допустимые координаты, а не один анзац.

\paragraph{Конкретная последовательность дальнейших задач.}
\begin{enumerate}
\item Для заданного $H(a,b,c)$ сначала применять точные чебышёвские
инволюции и сдвиги индексов; сохранять сам автоморфизм и обратный,
а не только конечную тройку индексов.
\item При кубической первой координате использовать матрицу остатков
в базисе $1,t,t^2$ и доказанный фильтр допустимых исключительных слоёв.
Для переноса с сохранённой второй координатой $C_7$ возможный
множитель имеет корни только в $x=0$. Естественный переход через
$x^2-2$ поэтому отбрасывается до большого символьного раскрытия.
\item Искать другой полином параметра в классе $R+I(\gamma)$.
Для кандидата проверять геометрические нули соответствующего
якобиева LND и строить слайс точным обращением по модулю полного
исключительного множителя. Отсутствие критических точек у кандидата
само по себе не удостоверяет координатность.
\item Для $H(4,5,7)$ начинать поиск координаты параметра со степени 4:
единственный полином степеней до 3 уже исключён. Для каждого найденного
полинома нужны компаньоны и полиномиальная обратная замена.
\item Для связанного кольца $B=\ker D_\gamma$ изучать существование
LND, перемещающего $t$. Дихотомия $\mathrm{ML}(B)\in\{k,k[t]\}$
уточняет цель, но не предоставляет алгоритма перечисления всех LND.
\end{enumerate}

\paragraph{Какие следствия уже можно использовать.}
Для удостоверенных моделей известны: два образующих идеала кривой
в координатах выпрямителя; свободное действие $\mathbb G_a$, одна
из орбит которого есть эта кривая; полиномиальные ядра связанных
пятимерных LND; явная тривиализация указанных конкретных действий;
минимальность степени числового поля определения, равной двум,
для нетривиальных вещественных базовых узлов. Ни одно из этих
утверждений не требует мотивных или гомологических методов.

\section{Материалы и повторение вычислений}

Архив содержит это досье с исходниками LaTeX, новые вычислительные
скрипты, JSON-сертификаты, журналы и манифест исходной версии репозитория.
Внешний репозиторий указан как фиксированная зависимость; его
полную копию архив не дублирует. Файл \texttt{bootstrap\_source.py}
получает только нужные файлы указанного коммита и проверяет Git-хеш
каждого файла. Это нужно лишь проверкам, зависящим от базового выпрямителя.
Самостоятельные сертификаты минимальных степеней слайсов используют
стандартную библиотеку Python и не требуют загрузки репозитория.

\begin{verbatim}
python -m pip install -r requirements.txt
python bootstrap_source.py
python research/repo_audit/independent_base_lnd.py
python research/harmonic_literature/periodic_shears.py
python research/lnd_certificate/verify_lnd.py
python research/lnd_certificate/verify_minimal_slice.py
python research/repo_audit/review_minimal_slice.py
python research/harmonic_compute/verify_new_results.py
python research/harmonic_compute/reduction_scan.py
\end{verbatim}
Подробный перечень команд дан в \texttt{README.txt}. Результаты
исходного запуска находятся рядом со скриптами. Новые проверки
используют явные исключения или явные проверки результата, где
это предусмотрено кодом; запускать команды следует обычным Python,
без отключения стандартных проверок флагом \texttt{-O}.

\paragraph{Основные первоисточники и зачем их читать.}
\begin{enumerate}
\item G.~Freudenburg, J.~Freudenburg,
\href{https://arxiv.org/abs/0902.3440}{Curves Defined by Chebyshev Polynomials}
(2009), предложения 3.1, 4.1, 4.2. Критерий комплексного вложения,
узлы плоской проекции и алгебраическая редукция индексов.
\item P.-V.~Koseleff, D.~Pecker,
\href{https://arxiv.org/abs/0909.3281}{Chebyshev diagrams for two-bridge knots};
\href{https://arxiv.org/abs/1203.4376}{Harmonic knots}.
Топология гармонических узлов, в частности классификация при $a\leq4$;
этот результат следует отделять от комплексной полиномиальной эквивалентности.
\item A.~van den Essen, P.~van Rossum,
\href{https://doi.org/10.1090/S0002-9939-01-06212-8}{Triangular derivations related to problems on affine $n$-space}
(2002), Theorem 4.1, Conjectures 5.1--5.3. Точное место гипотезы E
и связь между выпрямлением, ядрами LND и конкретным кандидатом
к проблеме сокращения.
\item D.~Lewis,
\href{https://arxiv.org/abs/2203.08178}{A unified approach to embeddings of a line in 3-space}
(2022), Lemma 1, Theorem 2, Corollary 4. Подъём параметра до
координаты и устранение полюсов рациональных преобразований.
\item K.~Masuda,
\href{https://doi.org/10.1016/j.jalgebra.2008.12.012}{Homogeneous locally nilpotent derivations having slices and embeddings of affine spaces}
(2009), и
\href{https://sci-tech.ksc.kwansei.ac.jp/~kayo/cmn/papers/eembed.pdf}{Equivariant embeddings and semi-invariant locally nilpotent derivations}.
Связь с градуировками и симметриями; дополнительные условия на веса
нужны для обратных импликаций.
\end{enumerate}
Нумерация и формулировки сверялись с текстами источников, а не только
с поисковыми аннотациями. Даты, автоматически вставленные в HTML
старых arXiv-статей, не использовались как даты публикации.
